\section{Comparing the methods}
\label{sec:comparison}
The three methods share query embedding, preparation, routing, filtering, and the optional
final reranker. Their difference is the work used to obtain candidates from the selected pool.
Keeping those shared costs visible lets us compare both the index work and the complete query.

\Needspace{12\baselineskip}
\subsection{Time per query}
Using the earlier definitions,
\begin{align}
T_A={}&\Tcommon+P c_{\rm range}
 +\Nelig Gc_{\rm lut}+H_s(\Nelig,R)c_{\rm heap}
 +T_{\rm rank}(R)+\Tfinish(R'),\label{eq:final-A}\\[3pt]
T_B={}&\Tcommon+T_{\rm order}+T_{\rm roots}
 +\sum_j s_jW_jc_{\rm split}+\sum_j\ell_jW_jc_{\rm leafword}\notag\\
 &+Vc_{\rm node}+FGc_{\rm lut,leaf}
 +H_s(F,R)c_{\rm heap}+T_{\rm rank}(R)+\Tfinish(R'),\label{eq:final-B}\\[3pt]
T_C={}&\Tcommon+T_{\rm keyprep}+E c_{\rm enum}+C^{\rm copy}c_{\rm copy}
 +\sum_j\left(H_jc_{\rm lookup}+F_jc_{\rm post}\right)\notag\\
 &+FGc_{\rm lut,posting}+H_s(F,R)c_{\rm heap}
 +T_{\rm rank}(R)+\Tfinish(R').\label{eq:final-C}
\end{align}
This C expression follows the implemented shared stream: $E$ counts its enumeration work once,
and $C^{\rm copy}$ counts its mask-copy work once. Cluster sums cover actual directory lookups
and postings. $F=\sum_jF_j$, and the current scorer evaluates all $G$ groups per candidate.
A tail-only or full-key score-reuse implementation would substitute its actual group count;
those savings are not credited to the measurements here.

\Needspace{12\baselineskip}
\subsection{Speeding up the whole query}
If two configurations have identical embedding and final-stage costs, write their shared time
as $T_0$. Replacing index time $t_A$ by $t_B$ gives speedup
\[
\frac{T_0+t_A}{T_0+t_B},
\]
not $t_A/t_B$. Even removing index time entirely cannot exceed $1+t_A/T_0$ when $T_0>0$.
If candidate counts or final-vector fetches differ, those costs are no longer shared and must
be evaluated separately. This is why cached-query index timings and complete text-query timings
should both be reported rather than substituted for one another.

\Needspace{12\baselineskip}
\subsection{Stored index memory}
The earlier layouts count the bytes needed for the data itself. The implemented search cores
also reserve space, so we count their allocated memory explicitly:
\begin{align}
M_{A,\rm core}&=M_{\rm codes,cap}+M_{\rm IDs,cap}+M_{\rm other,A},\label{eq:final-storageA}\\
M_{B,\rm core}&=M_{\rm codes,cap}+M_{\rm IDs,cap}+M_{\rm planes,cap}
 +M_{\rm other,B},\label{eq:final-storageB}\\
M_{C,\rm core}&=M_{\rm codes,cap}+M_{\rm IDs,cap}
 +\sum_j(U_jb_{\rm node}+L_jb_{\rm ptr})+M_{\rm other,C}.
\label{eq:final-storageC}
\end{align}
``cap'' means allocated array capacity. The remaining terms cover bucket objects, the outer
label map and allocator bookkeeping not already counted. C's IDs are its postings; there is
no second $N$-row posting permutation. Its hash-node and bucket-pointer terms use the separate
chaining layout in Equation~\eqref{eq:native-key-storage}.

The data stored in the planes takes $d b_w\sum_j\ceil{n_j/w}$; allocated capacity is at least that amount.
Storage sums run over all clusters, not just the probed ones. Add any retained routing index,
filters and fine-ranking vectors to the C++ index to obtain the index memory kept loaded during search. Input arrays
kept only during construction belong to the peak while building, not the memory kept during search.

\Needspace{12\baselineskip}
\subsection{Peak query memory}
\fig{memory}{Count stored index data, runtime reserve, and temporary memory used at the same time separately. Check when each stage allocates and releases its buffers.}{.85}

The temporary memory used during search is
\begin{align}
Q_{A,\rm search}&=\Qbase+Q_{\rm scan\ workspace},\\
Q_{B,\rm search}&=\Qbase+Q_{\rm order}
 +\max_t[\gamma_B Q_{B,\rm frontier}(t)+Q_{\rm current}(t)+Q_{\rm children}(t)],\\
Q_{C,\rm search}&=\Qbase+Q_{\rm keyorder}+Q_{\rm selected}
 +\max_t[\gamma_C Q_{\rm enum}(t)b_{\rm state}].
\end{align}
C has one queue for the common key coordinates, not one queue per probed cluster.
$Q_{\rm keyorder}$ stores the sorted coordinate indices; $Q_{\rm selected}$ stores the selected
bucket pointers. $\gamma_C Q_{\rm enum}$ accounts for queue capacity rather than only live entries.
The candidate heap is already in $\Qbase$, so it is not added again. The current C implementation
reads posting ranges directly and allocates no separate gathered-document buffer.

At the Python boundary, native result vectors and output arrays may coexist. Include that
interface workspace when evaluating the full request peak. Whole-process observations include
it even when individual interface allocations are not separately measured.

Small iterator state is included in workspace terms. Embedding and final reranking can have
more temporary memory than search, so use Equation~\eqref{eq:peak} to obtain the full query peak.
With loaded index size $M_{m,\rm resident}$ for method $m$, the RAM constraint is
\begin{equation}
\boxed{M_{m,\rm resident}+M_{\rm model}+M_{\rm runtime}
       +\max_t\sum_{u\in\mathcal A(t)}Q_{m,u}(t)\leq M_{\rm RAM}.}
\label{eq:final-ram}
\end{equation}
$M_{\rm model}$ is the shared memory for the loaded embedding model if encoding runs locally; it is zero in the memory being counted here when an external service performs encoding. It is not multiplied by the number of queries.

For a conservative bound, replace the sum by $JQ_{m,\rm peak,max}$. A loose upper bound that
exceeds RAM does not prove that a tighter implementation cannot fit. An observed peak below
the limit does not prove that every possible query fits. These are different claims.

\Needspace{12\baselineskip}
The native measurements use the allocation choices in Equation~\eqref{eq:native-key-storage}:
C reorders codes and IDs together, so its extra posting permutation is absent. Its directory
uses allocated nodes and bucket pointers. Use that layout when connecting the general formulas
to the measured results.

\subsection{The methods side by side}
\small
\begin{longtable}{@{}P{.16\linewidth}P{.25\linewidth}P{.25\linewidth}P{.25\linewidth}@{}}
\toprule & A: full scan & B: bitplanes & C: key enumeration\\\midrule\endfirsthead
\toprule & A: full scan & B: bitplanes & C: key enumeration\\\midrule\endhead
Additional index & none beyond common layout & $d b_w\sum_jW_j$ & key directory; IDs already counted\\
Main work & every eligible code & split masks, leaf masks, full scores, pending nodes & empty and occupied keys, enumeration, postings, full scores\\
Variable counts & $\Nelig$ & $s_j,\ell_j,f_j,V$ & $H_j,F_j,E,Q_{\rm enum}^{\max}$\\
Main memory concern & selection or batch score buffers & many full-width pending masks & pending flip sets and interface outputs\\
When local results are exact & scan all eligible rows & valid bounds, sufficient traversal, correct ties & full-key weighted order or a valid full-score bound\\
Settings that affect recall & routing probes; candidate/rerank boundary & node budget, visit policy, earlier routing & prefix width, candidate or key budget, earlier routing\\
What must be measured & scoring cost, routing quality & skipped scores versus bitmap/queue cost and recall & number of empty keys versus useful candidates, timing and recall\\\bottomrule
\end{longtable}
\normalsize

\Needspace{12\baselineskip}
\subsection{What changes time and memory use}
\textbf{Documents per physical cluster, $n_j$.} A grows roughly linearly with eligible row
count. In B, a larger cluster also widens every active bitmap. In C, more documents per key can
reduce the number of empty probes needed, while increasing posting and scoring work.

\textbf{Binary dimension, $d$.} Longer codes increase score terms and stored bits. B may gain
more from avoiding full scores, but also stores more planes and can traverse more coordinates.
For full-key C, the number of possible keys grows exponentially. For short-key C, $h$ and $d$
are separate parameters.

\textbf{Probes, $P$.} More clusters can recover missed neighbors. They also increase eligible
rows, root masks, active key streams and candidate competition. Cluster sizes and the query's
routing behavior matter, not just the number $P$.

\textbf{Leaf size and node budget.} A small leaf threshold can trade scoring for more mask work.
A finite node budget caps visits but does not specify how many useful leaves are reached.
The best choice depends on the query and data, not a general formula in which recall always moves in one direction.

\textbf{Lookup width and candidate target.} A short key has denser postings and a smaller key
space, but leaves more score information out. A longer key can improve selection of useful candidates
and simultaneously make empty-key enumeration costly. A strong prefix can improve recall;
that does not remove the enumeration cost.

\textbf{Scoring code and memory system.} SIMD, whether data stays in the CPU cache, batch size, alignment, layout and
thread count change the effective coefficients. Keep them fixed for an algorithm comparison,
or report their changes as a separate experiment.

\Needspace{12\baselineskip}
\subsection{Choosing settings within a memory limit}
Let $\theta_m$ contain the chosen method's tunable parameters and let $r_0$ be a recall target.
Choose which latency measure to minimize, for example median or p95 across queries. We can write that choice as
\begin{equation}
\theta_m^*\in\arg\min_{\theta_m}\operatorname{Quantile}_{p}[T_m(q;\theta_m)]
\end{equation}
subject to the RAM condition in Equation~\eqref{eq:final-ram} and
\begin{equation}
\E_q[\operatorname{Recall@}K(q;\theta_m)]\geq r_0.
\label{eq:optimization}
\end{equation}
The expectation specifies an average-recall target. A requirement that almost every query
reach $r_0$ is stronger and needs a different probability constraint.

\par\noindent\begin{minipage}{\linewidth}
\begin{lstlisting}[caption={A defensible finite-grid experiment.}]
ChooseConfiguration(corpus, queries, fixed_reference, RAM_limit, targets):
    build or cache every configuration's required index
    compute exact reference results once for each query
    on tuning queries:
        for every configuration in the declared grid:
            measure latency, quality, work counts and peak RAM
            keep failures and empty returns in the record
    choose qualifying configurations for each target
    on separate test queries:
        evaluate those choices without retuning
    report lowest tested latency only where the target and RAM checks pass
    distinguish measured results from untested model predictions
\end{lstlisting}
\end{minipage}\par

The parameters are often integers: cluster count, probes, lookup bits, and visited nodes. A
small stated set of combinations handles these choices directly. Differentiation can simplify a smooth
part of the cost model, but it does not determine whether the recall constraint is satisfied.
For example, a model $T(C)=aC+bNP/C$ has a stationary point, where its derivative is zero, at
$C=\sqrt{bNP/a}$ when $a,b>0$ and $P$ is held fixed. Check nearby integer values that meet the constraints, including
their memory use. If the probes required for the recall target change with $C$, that derivative no
longer solves the complete problem.

Choose a setting using queries for fitting and tuning, then check that fixed choice on
separate queries. Optimizing an inaccurate recall model can select the configuration where
that model makes its most optimistic error. At 99\% recall, a predicted 99\% is insufficient;
report actual query-level quality and its uncertainty. The existing normal prediction's
real-data errors are too large to establish such a winner.

If $d$ is also varied, keep the reference used to judge quality fixed across choices, for example the same
long-vector top $K$ or the same relevance labels. Recomputing a different binary ground truth
for every $d$ makes a tiny representation look perfectly accurate against its own changed task.

\subsection{Billion-document estimates}
\label{sec:billion}
The C++ cost models were checked through 8M rows. The bitmap test reaches the full width
of a 1B-row mask, but no complete 1B index was built or searched. The numbers below are
estimates under the same 256-bit score, top 100, one CPU thread and stated distribution of
strong coordinates. A different machine or pattern of memory reads can change the costs.

Start with memory, because it can reject a configuration without estimating its latency.
For the current layout at $N=10^9$ and $d=256$,
\[
M_{\rm codes}=32\ \mathrm{GB},\qquad M_{\rm IDs}=8\ \mathrm{GB}.
\]
A and C need at least 40 GB before directories, containers and query buffers. B retains a
second transposed copy, adding at least 32 GB. Therefore all three fail a 32 GB limit, and B
also fails a 64 GB limit. Changing probes does not remove stored documents. These rejections
apply to the current in-memory layouts; they are not claims about every compressed or disk-based design.

The planning figures below allow room for up to twice the stored row-ID count, a small
allowance for directories and containers, and a configurable 1 GB reserve for running queries.
That reserve is assumed, not measured. Payload counts exclude overhead: if even those bytes
exceed the budget, rejection is certain for this layout. A planned fit is only an estimate. Building an index from full input copies can require more RAM than serving it.

\begin{center}\small
\begin{tabular}{@{}lrrr@{}}\toprule
1B-row path & Payload lower bound & Planned query RAM & 64 GB budget\\\midrule
A, known strong-prefix routing & 40 GB & about 49 GB & Modeled fit\\
C, global 13-bit key & 40 GB & about 49 GB & Modeled fit\\
B, one global cluster & 72 GB & about 83 GB & Rejected\\
B, same prefix route, scan-like leaf & 72 GB & about 81 GB & Rejected\\\bottomrule
\end{tabular}\normalsize\end{center}
All these planning figures fit a 1 TB budget. The model does not claim that a 1 TB machine has
the M5 Max's exact memory behavior. The measured row buffers were much smaller than a 40 GB store.

\subsection{Fixed strong prefix}
With thirteen fixed strong coordinates, the ideal group contains an expected
$10^9/2^{13}\approx122{,}070$ rows. Under the independent-bit assumption, a binomial calculation
describes the chance of having fewer than 100. A Chernoff bound puts that probability below
$10^{-26464}$ in this example. This extremely small bound belongs to the stated distribution;
it is not an embedding-accuracy guarantee.

A can route directly to that group and scan it. C can look up the same key and score its posting
range. Their modeled costs are both roughly 3 ms. The data being read after routing is only a
small part of the corpus, although all codes and IDs are still stored.

B must get the same opportunity: it can use that prefix routing and choose a leaf large enough
to scan the selected group. An all-ones leaf-mask profile shows a similar scoring cost. Thus
B's scan fallback is also estimated at roughly 3 ms, while storing a much larger index. Forcing B to
traverse the entire billion-row bitmap would be an unfair parameter choice in this case.
The model is not precise enough to distinguish a timing gap of a few percent here.
A or C needs less memory, but the estimate does not prove one method is always faster.

\subsection{Changing strong coordinates}
A fixed prefix no longer supplies the same filter. The measured B path can select its columns
from the actual query. Keeping thirteen cuts and scaling the leaf threshold to 160,000 gives
an expected final set of about 122,070 rows. Its mask payload at the last split is
\[
(13+2)\ceil{10^9/64}\times8=1.875\ \mathrm{GB}.
\]
The model fitted to C++ runs estimates about 0.10 s for this global B path. If we double
the part of its cost that grows with corpus size, the estimate becomes about 0.20 s. This
checks the effect of a higher assumed cost; it is not a confidence interval or promised
bound. The model estimates about 25 s for a global scan.

That comparison does not establish that B beats every possible IVF configuration. Let $f$ be
the fraction of the corpus A actually scores, and let $T_{\rm route}$ be the time needed to
obtain those candidates. A's approximate cost is
\[
T_A\approx T_{\rm route}+fN c_{\rm scan}.
\]
Ignoring routing time gives A the most favorable comparison. At the current prices, its
scanned fraction must be below about 0.41\% to undercut the 0.10 s B point, or below about
0.82\% if B's variable cost doubles. It must still meet the recall target at that fraction.
If routing already produces a sufficiently small, accurate pool cheaply, scanning can win.
If it cannot, a global or less heavily clustered B path is a candidate worth testing.

\fig{billion-break-even}{A conditional break-even calculation at 1B rows with sufficient RAM. The scan line excludes routing time, favoring A. B is one fixed global operating point, not a retuned curve. Routing recall and the concentrated-query assumption must be checked separately.}{.94}

B can also use a router. In that case, substitute the physical cluster widths and visited
work into the earlier sum over clusters. The graph does not solve that larger optimization
problem or assume a relevant document is always in the probed clusters.

For C, the number of strong coordinates actually captured by its key matters. Under
Equation~\eqref{eq:captured-strong} and the stated bound conditions, a query capturing none
can still require scoring nearly all rows; capturing one leaves roughly half. A fixed short
key can therefore remain expensive even when it has many occupied entries and few empty
lookups. It is not enough to argue that finding a nonempty key is cheap.

\subsection{What the tests support}
The completed tests establish three different outcomes within their declared spaces:
\begin{itemize}
\item On the real Nomic sample, no alternative speedup was established. Even a one-bit key
control did not improve latency. Recall targets must retain their sampling uncertainty.
\item With a fixed strong prefix, direct routing and integrated key lookup are closely related
and both avoid most scoring. C had a small measured implementation advantage in the controlled runs.
\item With changing strong coordinates, B's query-dependent column order produced a large,
independently checked advantage in the constructed data. Increasing the leaf size with corpus
size preserved a useful path; keeping a tiny fixed leaf would not.
\end{itemize}
At a billion rows, the minimum memory counts are exact for this layout. Query times and
the best routing settings remain estimates under stated assumptions. The formulas give
settings to test and the conditions under which costs would be equal. A finite parameter
grid, or an embedding model name, cannot identify the best setting for every workload.

The current calculation is \path{experiments/project_costs.py}; its inputs and all scenarios
are saved in \path{results/projected-cases-2026-09-11/}. The earlier
\path{reports/search-math/numerical-cases/calculate.py} is preserved as a historical assumed-cost
model. Its extra-posting and hash-slot layout must not be substituted for the measured native C layout.


\begin{resultbox}{What the comparison tells us}
A pays to score the selected pool. B pays to discover which scores can be avoided, and dense
bitmaps can make that discovery expensive. C pays to discover which keys contain useful
candidates, including empty keys and queue work. Each can be exact under a suitable stopping
rule. Which is fastest at a recall target is determined by the data, the representation,
the visited work, the memory layout and the hardware. The exact small examples explain how
recall can be calculated. The measured prediction errors show where the assumptions need more
evidence before they can guide a real configuration choice.
\end{resultbox}
